\section{Radial energies and the radiation potential}\label{sec:radial}

Throughout the proof the parameters \(n,\beta,\gamma,R\) are fixed as in
Theorem~\ref{thm:main}. We write
\[
 S=n\omega_n,\qquad A(r)=r^{n-1},\qquad Y=\Theta(1).
\]
Thus \(S A(r)\) is the area of the sphere of radius \(r\).

The interior trace on a bounded Lipschitz open set is a bounded linear
map from \(BV\) into boundary \(L^1\). With outward normal \(\nu\),
Gauss--Green reads
\[
 \int_\Omega f\,\operatorname{div}X+\int_\Omega X\cdot\dd Df
 =\int_{\partial\Omega}\Tr f\,X\cdot\nu\,\dd\mathcal H^{n-1}.
\]
This is the inner-normal formula of \cite[Theorem~2.9]{LV} with its sign
converted. We use the absolute-mean characterization
\[
 \lim_{\rho\downarrow0}\frac1{|B_\rho(x)\cap\Omega|}
 \int_{B_\rho(x)\cap\Omega}|f-\Tr f(x)|\dd y=0
 \quad\text{for }\mathcal H^{n-1}\text{-a.e. }x\in\partial\Omega
\]
from \cite[Corollary~1.29]{Gorny}.\footnote{The proof of this
corollary is omitted in the lecture notes, whose publication date is
not specified. The point in \cite[(2.5)]{LV} is printed as \(x\in\Omega\),
whereas \cite[(2.6)]{LV} is written on \(\partial\Omega\). The
absolute-mean statement used here is the boundary statement in
\cite[Corollary~1.29]{Gorny}.}
If a connected domain is required in that characterization, apply it
in one-sided Lipschitz chart cylinders. On smaller boundary patches the
small-ball intersections agree with those in \(\Omega\), and uniqueness
of the local trace gives the stated formula without connectedness of
\(\Omega\).


We first justify the minimum in the definition of the energy. An admissible
minimizing sequence is bounded in \(H^1(\Omega)\), since its values lie in
\([0,1]\). Pass to a weakly convergent subsequence. Tests by indicators show
that the limit still lies in \([0,1]\) and equals 1 almost everywhere on \(K\).
The bounded trace map gives
\[
 \|\Tr f\|_{L^1(\partial\Omega)}
 \le C_\Omega\bigl(\|f\|_{L^1(\Omega)}+\|\nabla f\|_{L^1(\Omega)}\bigr).
\]
Consequently every pairing of the trace with a bounded boundary function
is a continuous linear functional on \(H^1(\Omega)\). Traces of admissible
temperatures also lie in \([0,1]\). Since
\(\Theta''(s)=4\beta(s+\gamma)^3\ge0\), the supporting-line inequality at
the limiting trace proves weak lower semicontinuity of the boundary
integral. The gradient integral has the same property. The constant
temperature 1 is an admissible finite-energy test, so the minimum is attained.

\begin{lemma}\label{lem:radial}
For \(r>1\), let \(I_n(r)=\int_1^r \rho^{1-n}\dd\rho\). Then
\begin{equation}\label{eq:radial}
 \frac{E_\Theta(B_1,B_r)}{S}
 =\min_{0\le t\le1}\left\{\frac{(1-t)^2}{I_n(r)}+
                                      r^{n-1}\Theta(t)\right\}.
\end{equation}
The energy is continuous on \([1,R]\) when its value at 1 is defined
by the bare ball, \(E_\Theta(B_1,B_1)=S Y\). In particular its minimum
on this interval is attained.
\end{lemma}

\begin{proof}
For an admissible \(v\) on \(B_r\), let
\[
 m(\rho)=\frac1S\int_{\partial B_1}v(\rho\theta)
                                      \dd\mathcal H^{n-1}(\theta).
\]
Radial tests in the weak-derivative identity give
\(m'(\rho)=S^{-1}\int_{\partial B_1}\nabla v(\rho\theta)\cdot\theta
\dd\mathcal H^{n-1}(\theta)\). Cauchy--Schwarz therefore yields
\[
 S\int_1^r A(\rho)|m'(\rho)|^2\dd\rho
 \le\int_{B_r\setminus B_1}|\nabla v|^2\dd x.
\]
In particular \(m\) has an absolutely continuous representative on
\([1,r]\). Applying the same argument across radius 1 shows \(m(1)=1\).
Its value \(m(r)\) is the mean of the actual boundary trace. To see this
without selecting a boundary slice, use Gauss--Green with the radial
field \(f(\rho)\theta\), zero near the origin. Polar integration turns
its left side into
\[
 S\int_0^r\bigl[m(Af)'+m'Af\bigr]\dd\rho
       =S A(r)f(r)m(r),
\]
whereas its boundary side is
\(f(r)\int_{\partial B_r}\Tr v\,\dd\mathcal H^{n-1}\).
Here \(m'=0\) in the core.

Put \(t=m(r)\in[0,1]\). Weighted Cauchy--Schwarz gives
\((1-t)^2\le I_n(r)\int_1^r A|m'|^2\).
The supporting line for \(\Theta\) at the boundary mean gives
\(\int_{\partial B_r}\Theta(\Tr v)\ge S A(r)\Theta(t)\).
Conversely the radial temperature
\[
 u(\rho)=1-(1-t)\frac{I_n(\rho)}{I_n(r)},\qquad 1\le\rho\le r,
\]
extended by 1 in the core, has exactly the cost on the right of
\eqref{eq:radial}. This proves the formula.

Compact minimization in \(t\) gives continuity for \(r>1\).
For a minimizing \(t_r\), the test \(t=1\) implies
\[
 (1-t_r)^2\le I_n(r)r^{n-1}Y\longrightarrow0\qquad(r\downarrow1).
\]
Thus its boundary term tends to \(Y\), while the entire expression
is at most \(r^{n-1}Y\to Y\). This proves continuity at 1 and the
assertion about the minimum. The integral definition includes
\(I_2(r)=\log r\), so no separate singular formula occurs in dimension two.
\end{proof}

We next identify the scalar property used to remove jumps. The potential
\(\Theta\) is strictly increasing on \([0,1]\), because
\(\Theta'(s)>0\) for \(s>0\). Let \(Q:[0,Y]\to[0,1]\) be its continuous
inverse, and define
\begin{equation}\label{eq:GH}
 G(y)=\frac{\Theta'(Q(y))^2}{4},\qquad
 H(y)=\frac{y^2}{G(y)}\quad(y>0),\qquad G(0)=H(0)=0.
\end{equation}
In potential variables \(y=\Theta(u)\), the Dirichlet density is
\(|y'|^2/(4G(y))\). The function \(H\) appears when this potential
is multiplied by a constant.

\begin{lemma}\label{lem:concavity}
The function \(G\) is Lipschitz on \([0,Y]\), positive on \((0,Y]\),
and \(H\) is concave on \([0,Y]\).
\end{lemma}

\begin{proof}
For \(y>0\), inverse differentiation gives
\[
 G'(y)=\frac{\Theta''(Q(y))}{2}
       =2\beta(Q(y)+\gamma)^3.
\]
This is bounded by \(2\beta(1+\gamma)^3\), and \(G\) is continuous at
zero. The Lipschitz assertion follows.
If \(\gamma=0\), then
\[
 H(y)=\frac4{25}\left(\frac{5y}{\beta}\right)^{2/5},
\]
which has the required concavity. Suppose \(\gamma>0\), put \(z=s/\gamma\),
and set
\[
 P(z)=z^5/5+z^4+2z^3+2z^2.
\]
Then \(\Theta(s)=\beta\gamma^5P(z)\) and
\(H(\Theta(s))=4\gamma^2P(z)^2/P'(z)^2\).
For \(h=4P^2/P'^2\), differentiation gives
\[
 h'=8P/P'-8P^2P''/P'^3,\qquad
 h''=8-24PP''/P'^2-8P^2P'''/P'^3+24P^2P''^2/P'^4,
\]
and hence
\begin{equation}\label{eq:concavitynumerator}
 h''P'-h'P''=-\frac{8}{P'^3}
          \bigl[P^2P'P'''-(P'^2-2PP'')^2\bigr].
\end{equation}
The second derivative of \(h\) with respect to \(P\) is the left
side of \eqref{eq:concavitynumerator} divided by \(P'^3\).
For \(q=5P=z^5+5z^4+10z^3+10z^2\), direct multiplication gives
\begin{align}
 q^2q'q'''-(q'^2-2qq'')^2
 =25z^5\bigl(&3z^{11}+48z^{10}+360z^9+1680z^8+5508z^7 \notag\\
 &+13584z^6+26136z^5+39600z^4+46080z^3 \notag\\
 &+38400z^2+20000z+4800\bigr). \label{eq:positivepoly}
\end{align}
All coefficients are positive. The expression in brackets in
\eqref{eq:concavitynumerator} is \(1/625\) times the left side of
\eqref{eq:positivepoly}. Since \(P'>0\), this proves strict concavity
at positive potential values; all changes of scale have positive factors.
Finally \(P(z)\sim2z^2\), \(P'(z)\sim4z\), so
\(H(\Theta(s))\sim s^2\to0\) as \(s\downarrow0\).
The continuous zero extension preserves concavity.
\end{proof}
